
\documentclass[11pt]{article}
\usepackage[a4paper,margin=1in]{geometry}
\usepackage{amsmath,amssymb,hyperref}
\title{Counting Addendum for the de Bruijn Sheaf/Homology Framework}
\author{Generated addendum}
\date{\today}
\begin{document}
\maketitle

\section{Counting de Bruijn cycles}

The number of de Bruijn cycles of order \(n\) over an alphabet of size \(q\) is
\[
N(q,n)=\frac{(q!)^{q^{n-1}}}{q^n}.
\]
For \(q=2\),
\[
N(2,n)=2^{2^{n-1}-n}.
\]
This follows by applying the BEST theorem to \(B(q,n-1)\).

\section{Homology}

An Eulerian circuit is a full-support integral 1-cycle:
\[
\partial c=0.
\]
Thus \(N(q,n)\) counts full-coverage global cyclic sections of the word sheaf.

\section{Pattern-level decomposition}

Let \(\overline B(q,n)=B(q,n)/S_q\). Then
\[
N(q,n)=\sum_S |\operatorname{Lift}(S)|,
\]
where \(S\) ranges over pattern skeletons in the quotient graph.

\section{Constrained systems}

For constraints \(\mathcal C\),
\[
N_{\mathcal C}(q,n)
=
\sum_{\text{valid }S}
|\operatorname{Lift}_{\mathcal C}(S)|.
\]
Near-\(G\) defects appear as nonzero boundaries:
\[
\partial c\neq0.
\]

\section{Applications}

DNA/RNA reconstruction, evidence-based prescription, and weather/atmospheric disturbance analysis all fit this form: local sections plus overlap restrictions yield global reconstruction problems, with defects measured by boundaries.

\end{document}
